\section{Introduction}
\label{sec:intro}

Coccidiosis alone was estimated to cost the global chicken industry about \pounds10.4 billion at 2016
prices~\cite{blake2020}. Coccidiosis, salmonellosis and Newcastle disease change the colour and consistency of
droppings, and farmers already look at droppings to judge flock health. This has made dropping photographs a popular
target for deep learning. On the public Tanzanian collection of Machuve et al.~\cite{machuve2022,machuve2021farm},
convolutional classifiers reach 98 to 99\% accuracy~\cite{machuve2022,dhungana2025,degu2023,reyad2024,huda2024}, and
several groups have proposed phone, web or edge systems around them~\cite{degu2023,sakib2026,dhungana2025}.

These accuracies come from random splits of a single collection. A random split by photograph places pictures from
the same farm visit, and sometimes the same picture, on both sides of the split~\cite{barz2020,kapoor2023}. The
original authors wrote that they did not expect their model to be affected by distribution shift~\cite{machuve2022}.
This has not been tested for the four-class task. The same group released a second collection whose labels were
confirmed by PCR~\cite{machuve2021pcr}, which offers a direct test. A second gap is the camera. Published systems
classify close-up photographs taken by a person. A camera fixed above a pen sees each dropping at 8 to 30 pixels,
among birds, feeders and litter, and nobody is there to restart the device when software fails.

We address both gaps with one detector and one honest protocol. Our contributions are as follows.

\begin{enumerate}
\item A benchmark of the classifiers earlier work used (MobileNetV2, EfficientNet-B0) and a ResNet-50 reference,
  retrained on the same split as our YOLO11n detector and tested on the same 679 farm photographs and on 1,255
  lab-confirmed photographs that no model saw. On the farm test our detector calls 97.9\% of photographs correctly,
  in line with published numbers. On the lab-confirmed photographs it calls 67.1\% correctly
  \todo{baseline numbers}. To our knowledge this is the first such test for the four-class task.
\item An audit of the split. We find that 49 of the 679 test photographs (7.2\%) are exact copies of training
  photographs, and we show that removing them changes the farm accuracy by only 0.1 points. We also show that the
  detector's box score is not a calibrated probability.
\item A pen-scale fine-tune trained with real droppings pasted onto barn video frames, split by camera. On cameras it
  never saw, recall rises from 0.29 to 0.80 and precision from 0.11 to 0.93, at a cost of 0.02 mAP50 on close-ups.
\item A complete pipeline running on a Raspberry Pi 5 under the QNX microkernel, with every model timed on the device.
  Each stage is an isolated process. When we kill any stage, it is restarted in 20 to 33\,ms while the others keep
  running. To our knowledge no peer-reviewed work has run poultry health vision on a microkernel real-time operating
  system.
\end{enumerate}

We release the code, the per-photograph predictions, the duplicate hashes and every figure with the table behind it.
